% SGA 3，第 VII_B 篇，英文重建稿——第 1 轮。
% 范围：权威底本分篇第 42--43 页，第 1.6.1 节、其末尾的命题及注 1.6.2 全文，止于分篇第 43 页第 2 节之前。
% 合并阅读本页码： 555--556.
% 再版印刷页码： 544--545.
% 本范围内原版页边定位： 511.
% 本范围内编者注： 74--75.
% 权威底本： Polo--Gille Exp7B-13oct24.pdf,
% SHA-256 D303690E0195654BBE7C3DAA5D89EA23E24CD709A3575217234AC19AE8836694.
% 曾参考 Jacob Reinhold 的 CC BY 4.0 Markdown 作为比对谱系。
% 原英文 Loop 1 阶段说明：三幅图以直接取自 PDF 的精确源裁片保留，普通数学均为原生 LaTeX。当前中文工程已采用原生图。

\subsubsection*{1.6.1}
\phantomsection
\label{sga3:VIIB:1.6.1}

设 \(Y\) 为 平展 形式 \(k\)-簇，\(f:X\to Y\) 为 \(\operatorname{Vaf}/k\) 中的态射。则下列方块为笛卡尔，其中 \(\Gamma_f\) 是以 \(\operatorname{id}_X\) 和 \(f\) 为分量的图态射 \(X\to X\times Y\)：

\begin{center}
  \phantomsection
  \label{sga3:VIIB:1.6.1:cartesian-square}
  % SGA3-VII-DIAG-096
  % SGA3-VII-DIAG-096
% Authority: Exp7B-13oct24.pdf, local p. 42.
% Crop: figures/Exp7B_localp042_section161_cartesian_square.png
% Locator: tex/components/35_expose_VIIB_sections161_162_flat_etale_reflection_en.tex, Loop 1 line 27.
\begin{tikzcd}[row sep=2.8em,column sep=3.7em]
X \arrow[r,"\Gamma_f"] \arrow[d,"f"']
  & X\times Y \arrow[d,"{f\times\operatorname{id}_Y}"] \\
Y \arrow[r,"\Delta_Y"]
  & Y\times Y\, .
\end{tikzcd}

\end{center}

由此可知 \(\Gamma_f\) 为局部同构，因此若 \(\operatorname{pr}_Y\) 拓扑平坦，则 \[
  f=\operatorname{pr}_Y\circ\Gamma_f
\] 也拓扑平坦；例如当 \(X\) 在 \(k\) 上拓扑平坦时即如此。

反之，由于 \(Y\to k\) 拓扑平坦，若 \(f\) 拓扑平坦，则 \(X\to k\) 也拓扑平坦 \textup{（参见 \SGARefLink{sga3:VIIB:1.3.3}{1.3.3}）}。特别地，取 \(f\) 为 \SGARefLink{sga3:VIIB:1.6}{1.6} 中的典范态射 \[
  p:X\longrightarrow X_e
\]，便得到下述结论。

\medskip\noindent\textbf{命题。}\phantomsection \label{sga3:VIIB:1.6.1.prop}{\itshape 设 \(X\) 为 \(k\) 上的形式簇。态射 \(X\to X_e\) 拓扑平坦，当且仅当 \(X\) 在 \(k\) 上拓扑平坦。}

\medskip\noindent\textbf{注 1.6.2。}\footnote{\label{sga3:VIIB:nde:74}编者将此注插在这里；原文中它位于 \SGARefLink{sga3:VIIB:2.5.2:proposition}{2.5.2}。}\phantomsection \label{sga3:VIIB:1.6.2} 当 \(k\) 为\emph{完美}域时，函子 \[
% 权威底本分篇页 43; 合并阅读本页 556; 印刷页 545.
  X\longmapsto X_e
\] 也是包含函子 \[
  \operatorname{Vaf}^{\mathrm{\acute{e}t}}/k
  \lhook\joinrel\longrightarrow
  \operatorname{Vaf}/k.
\] 的右伴随。事实上，此时对每个 \(x\in X\) 有 \[
  \mathcal O_{X_e,x}=\kappa(x)
\]，而典范投影 \[
  \mathcal O_{X,x}\longrightarrow\kappa(x)
\] 定义态射 \(s:X_e\to X\)，它是 \(p:X\to X_e\) 的截面。对每个态射 \(f:X\to Y\)，下列图均交换：

\begin{center}
  \phantomsection
  \label{sga3:VIIB:1.6.2:local-ring-square}
  % SGA3-VII-DIAG-097
  % SGA3-VII-DIAG-097
% Authority: Exp7B-13oct24.pdf, local p. 43.
% Crop: figures/Exp7B_localp043_remark162_local_ring_square.png
% Locator: tex/components/35_expose_VIIB_sections161_162_flat_etale_reflection_en.tex, Loop 1 line 86.
\begin{tikzcd}[row sep=2.6em,column sep=3.5em]
\mathcal O_{Y,f(x)} \arrow[r] \arrow[d]
  & \mathcal O_{X,x} \arrow[d] \\
\kappa(f(x)) \arrow[r]
  & \kappa(x)
\end{tikzcd}

  \hspace{.07\textwidth}
  \phantomsection
  \label{sga3:VIIB:1.6.2:functoriality-square}
  % SGA3-VII-DIAG-098
  % SGA3-VII-DIAG-098
% Authority: Exp7B-13oct24.pdf, local p. 43.
% Crop: figures/Exp7B_localp043_remark162_functoriality_square.png
% Locator: tex/components/35_expose_VIIB_sections161_162_flat_etale_reflection_en.tex, Loop 1 line 91.
\begin{tikzcd}[row sep=2.6em,column sep=3.2em]
X \arrow[r,"f"] & Y \\
X_e \arrow[r,"f_e"] \arrow[u,"s_X"]
  & Y_e \arrow[u,"s_Y"']
\end{tikzcd}

\end{center}

因此 \(s\) 对 \(X\) 具有函子性，由此可知 \(X\mapsto X_e\) 是包含函子 \[
  \operatorname{Vaf}^{\mathrm{\acute{e}t}}/k
  \lhook\joinrel\longrightarrow
  \operatorname{Vaf}/k.
\] 的右伴随。因而，作为右伴随，函子 \(X\mapsto X_e\) 与 \(\operatorname{Vaf}^{\mathrm{\acute{e}t}}/k\) 中存在的投射极限交换，\footnote{\label{sga3:VIIB:nde:75}有限投射极限即存在。}特别地，与有限乘积交换。

\noindent（也可以直接验证：对仿射代数为 \(A\) 的每个形式 \(k\)-簇 \(X\)，\(X_e\) 的仿射代数是 \(A\) 对其根 \(r(A)\) 的商；由于 \(k\) 完美，\(\mathcal O_{X,x}\widehat{\otimes}\mathcal O_{Y,y}\) 对其根的商就是代数 \(\kappa(x)\otimes_k\kappa(y)\)，因为后者半单。）


% 下一游标： Exp7B 分篇页 43, 合并阅读本页 556, 再版印刷页 545, 紧接第 2 节之前。
